\section{本讲小结}

理想导线确定节点，元件的两端连接决定串并联。节点方程来自电荷守恒，回路方程来自稳恒场中的电势变化相加；参考电流与绕行方向一旦选定，就应保持到计算结束。负电流说明实际方向相反，断路则只使支路电流为零，不能同时删去其中的电源升势。

节点法以电势为未知量，网孔法以循环电流为未知量，共用电阻中的实际电流须按贡献相减。理想电压源跨未知节点时，超节点只作为共同守恒区域，两端仍保持规定的电势差。理想电流源规定电流，其电压由外网络决定。

线性网络可叠加各独立电源产生的电压和电流，关闭电动势时短路理想电压源、关闭理想电流源时开路，并保留原有电阻。功率包含电流的平方及交叉项，应在叠加之后重新计算。温度使电阻改变时，还须重新检查线性前提。

两端等效保持外部端电压与输出电流关系，开路电压给出截距，看入电阻给出斜率。等效时负载留在网络之外，接回负载后再回算原网络内部各支路。桥平衡由两中点电势相等确定，不平衡时则可把桥臂作为负载求解。全网功率核验应使用原网络全部电源与耗散元件，不能用等效网络内部的发热替代。
